\section{The model}
\label{sec:model}

\subsection{The table}

\shipped{} An organization's graph is one SQLite table in a file opened per
(organization, project), so that another tenant's rows are not in the database
being read~\cite{hip1198,sqlitedocs}. The project that selects the file comes
from a claim the identity service mints, never from the request body. A row is
an assertion.

\begin{definition}[Assertion]
An assertion is
$f = (e, r, v, n, \mathit{at}, \until, \mathit{seen}, \mathit{src},
\mathit{ev}, \mathit{by}, \mathit{conf})$ together with fields the server sets:
the write instant $\mathit{wrote}$, the knowable instant, a sequence number and
a hold bit. The entity $e$ is what is described, $r$ the relation and $v$ the
value; $n \in \{0,1\}$ says whether $v$ names another entity. With $n = 1$ the
assertion is an \emph{edge}, with $n = 0$ a \emph{property}. The interval
$[\mathit{at}, \until)$ is when the filer says the thing was so, with $\until$
absent for an open interval. $\mathit{seen}$ is when the filer says the
assertion became available, $\mathit{src}$ names who asserted it, $\mathit{ev}$
points at the record it came from, $\mathit{by}$ is the filing identity and
$\mathit{conf} \in [0,1]$ a confidence.
\end{definition}

An entity is not a row and is never created: it exists because something was
asserted about it. An edge and a property are the same kind of row, and $n$ is
the bit that tells a traversal which rows it may follow. The filing identity is
stamped from the validated principal and never read from the body. Instants are
stored as whole seconds.

\begin{definition}[Content address]
$\mathrm{id}(f) = H(e, r, v, n, \mathit{at}, \mathit{seen}, \mathit{src},
\mathit{ev}, \mathit{by}, \mathit{conf})$, with $\until$ appended when it is
set, where $H$ is SHA-256 over the length-prefixed fields. The write instant,
the knowable instant, the sequence number and the hold are excluded.
\end{definition}

The table has a uniqueness constraint on the content address and inserts with
\texttt{INSERT OR IGNORE}, so a redelivered assertion appends nothing and the
write reports it as a duplicate. Including the write instant would make every
retry a new row. An open interval keeps the address it had before $\until$
existed, so rows filed then still deduplicate. Two identities asserting the same
thing produce two rows, because $\mathit{by}$ is part of the address.

\subsection{The knowable instant}

\begin{definition}[Knowable]
$\knowable(f) = \max(\mathit{seen}(f), \mathit{wrote}(f))$.
\end{definition}

\shipped{} Admission requires a source and $\mathit{conf} \in [0,1]$, bounds the
size of the entity, relation, value and evidence, requires a value on an edge,
requires $\until > \mathit{at}$ when $\until$ is set, sets $\mathit{seen}$ to
$\mathit{at}$ when absent, requires $\mathit{seen} \ge \mathit{at}$, and refuses
an $\mathit{at}$ or a $\mathit{seen}$ more than five minutes ahead of the server
clock. An assertion dated further ahead would not become visible until that
instant and would change reads when it did. The end $\until$ is not bounded
by the clock: a term that is known to end next year may say so.

\subsection{Statements and versions}

\begin{definition}[Statement, version, retraction]
The \emph{statement} of an assertion is $s(f) = (e, r, v, n, \mathit{at})$. The
assertions sharing a statement are its \emph{versions}; they may differ in
$\until$ and in provenance. A \emph{retraction} is a property with an empty
value, $n = 0$ and $v = \varepsilon$.
\end{definition}

A version that changes $\until$ is how the end of a term is corrected without
deleting the first account. A retraction is how a claim is withdrawn: the pair
then reads as withdrawn, with the withdrawn statement and the retraction both
still readable.

\begin{definition}[Order]
$f \succ g$ when $\knowable(f) > \knowable(g)$; on equal instants, when
$\mathit{conf}(f) > \mathit{conf}(g)$; otherwise when
$\mathrm{id}(f) < \mathrm{id}(g)$ lexicographically.
\end{definition}

The order in \texttt{claim} has a first term, a rank the caller assigns to
sources. This store assigns every source the same rank, because ranking one
publisher above another is a claim about a domain for which it holds no
vocabulary, so the order reduces to the three terms above. Confidence
decides only between assertions known at the same instant.

\begin{definition}[Current versions]
For a transaction instant $\tk$, let $K_{\tk} = \{f : \knowable(f) \le \tk\}$.
For each statement $s$, the current versions $C_s(\tk)$ are the members of
$K_{\tk}$ with statement $s$ whose knowable instant is the latest among them.
All members sharing that latest instant are kept.
\end{definition}

A version therefore replaces another only by being learned later. Two versions
learned at the same instant stand together: in the CronQuestions data they are
two terms of one holder that begin in the same year, filed from two lines of the
source, and choosing one of them would discard a term (\S\ref{sec:eval-first}).

\subsection{Cardinality}

A relation holds one value at a time or several. An organization declares which
with an assertion (\texttt{relation:}$r$, \texttt{cardinality},
\texttt{one}\,$|$\,\texttt{many}), resolved by the same rule as any other pair
(\S\ref{sec:schema}). Undeclared, a relation holds one value when its pair holds
properties and several when it holds edges: a property has a value, and an
entity has neighbours. A position with one holder per term, read from the
position's side, is several values over time and at most a few at once; an
owner is one.

\subsection{Resolution}

\begin{definition}[Point, open]
A \emph{point} is a pair $p = (\tv, \tk)$ of a valid instant and a transaction
instant. A statement is \emph{open} at $\tv$ when $\mathit{at} \le \tv$ and either
$\until$ is absent or $\tv < \until$.
\end{definition}

Let $S$ be the union of $C_s(\tk)$ over the statements of one pair, restricted
to $\mathit{at} \le \tv$.

\begin{definition}[Resolution, one value]
Let $T \subseteq S$ be the members with the latest $\mathit{at}$, and $w$ the
$\succ$-greatest member of $T$. The pair \emph{holds} $w$ at $p$ when $w$ is not
a retraction and is open at $\tv$; otherwise it holds nothing. The
\emph{conflicts} are the other members of $T$ that are open, are not
retractions and name a different value, one per value.
\end{definition}

A later statement therefore ends an earlier one, a retraction ends what precedes
it, and a statement past its own $\until$ leaves the pair holding nothing rather
than falling back to an older value. Accounts of the same instant that disagree
remain visible as conflicts, and the response marks the pair contested.

\begin{definition}[Resolution, several values]
The pair holds every non-retraction $f \in S$ that is open at $\tv$ and is not
ended by a retraction $g \in S$, where $g$ ends $f$ when
$\mathit{at}(g) > \mathit{at}(f)$, or $\mathit{at}(g) = \mathit{at}(f)$ and
$g \succ f$.
\end{definition}

An edge is not ended by another edge, so a pair holds as many edges as were
asserted and not ended. A statement's own $\until$ ends only that statement; a
retraction ends every statement of its pair begun before it. The answer lists
what is held with the most recently begun first; that first statement is
reported as the winner.

\begin{remark}[What the previous version did]
Before \texttt{1fffdf408} a statement had no $\until$, every pair held one value,
and a read took one instant that bounded the knowable instant. A term was filed
as an assertion at its start and a retraction at its end. Under that rule the
retraction of one holder's term was the most recent statement about the
position, and it ended every other holder of the position as well
(\S\ref{sec:eval-first}).
\end{remark}

\subsection{Properties}

\begin{proposition}[Answers as known at an instant are stable]
\label{prop:stable}
Let an assertion be written at server instant $W$, and let $p = (\tv, \tk)$ with
$\tk$ in a second before that of $W$. Then the resolution of every pair at $p$,
and every declared cardinality read at $\tk$, is the same before and after the
write.
\end{proposition}
\begin{proof}
$\knowable(f) \ge \mathit{wrote}(f) = W > \tk$ at one-second resolution, so
$f \notin K_{\tk}$. Resolution at $p$ reads only $K_{\tk}$, through the pair's
rows and the cardinality declaration's rows, and both are filtered by
$\knowable \le \tk$ before any row limit is applied.
\end{proof}

\begin{proposition}[No backdating]
\label{prop:noback}
For every assertion $f$ and every $\tk < \mathit{wrote}(f)$, $f \notin K_{\tk}$,
whatever $\mathit{seen}(f)$ the filer supplied.
\end{proposition}
\begin{proof}
$\knowable(f) = \max(\mathit{seen}(f), \mathit{wrote}(f)) \ge
\mathit{wrote}(f) > \tk$.
\end{proof}

This is the guard HIP-1198 requires~\cite{hip1198}: a bound on a
caller-supplied instant is only as strong as the caller's honesty, so the bound
is on the derived instant. It constrains $\tk$ only. A caller may ask about any
valid instant, which is what the previous version did not allow.

\begin{proposition}[Resolution does not depend on storage order]
\label{prop:total}
$\succ$ is a strict total order on assertions with distinct content addresses.
The rows a pair holds at $p$ are therefore a function of the set of its rows in
$K_{\tk}$, not of the order in which they are read.
\end{proposition}
\begin{proof}
The first two terms compare totally ordered values; the last compares distinct
strings. Current versions are defined as a set. The one-value rule takes a
maximum of a finite set under a total order, which is unique, and the
several-value rule is a filter whose test uses only $\succ$ and the fields of
two rows.
\end{proof}

Section~\ref{sec:time} adds a fourth property, the condition under which a
valid-time answer does not depend on any knowable instant, because it is the
property the CronQuestions measurement tests.

\begin{remark}[Exceptions to Proposition~\ref{prop:stable}]
Erasure removes rows and can change an answer at any point; that is its purpose
(\S\ref{sec:erasure}). A read of one pair is bounded at 10{,}000 rows: past
that, resolution uses the rows begun most recently, and \op{resolve} reports
the truncation. The bound is applied after the filter on $\knowable$, so a
truncated answer is still stable in the sense of the proposition, but it may
differ from the answer the full pair would give.
\end{remark}

\subsection{The surface}

\shipped{} Table~\ref{tab:surface} lists the operations at \texttt{1fffdf408}.
Each except GraphQL is typed: one method, one path, one input type and one
output type, from which the published document, the generated clients and the
agent tool list are projected. The organization is never an input.

\begin{table}[h]
\centering
\small
\begin{tabular}{@{}lp{0.60\textwidth}@{}}
\toprule
route & what it answers \\
\midrule
\texttt{POST /v1/graph} & assert a batch; each member recorded, counted as a duplicate, or refused with a reason \\
\texttt{GET /v1/graph} & the rows matching an entity, relation or value, optionally bounded by \texttt{as\_of} and \texttt{as\_known}; resolves nothing \\
\texttt{GET /v1/graph/search} & the same read narrowed by a full-text match \\
\texttt{POST /v1/graph/resolve} & what one pair holds at a point, with conflicts \\
\texttt{POST /v1/graph/neighbors} & a bounded walk over edges in force at a point \\
\texttt{POST /v1/graph/path} & the shortest chain of edges in force between two entities \\
\texttt{POST /v1/graph/diff} & what began, was superseded and was retracted between two points \\
\texttt{GET /v1/graph/vocabulary} & relations in use, the declared schema, and the terms of the order \\
\texttt{POST /v1/graph/extract}, \texttt{/ingest} & the relations a document states, without and with recording them \\
\texttt{POST /v1/graph/communities} & a Louvain partition of the edges in force at a point \\
\texttt{POST /v1/graph/answer} & a map-reduce answer over the communities, with citations \\
\texttt{POST /v1/graph/erase} & erasure of one subject, with a receipt \\
\texttt{POST /v1/graph/graphql} & GraphQL whose resolvers call the operations above \\
\bottomrule
\end{tabular}
\caption{The operations of \texttt{/v1/graph} at \texttt{1fffdf408}. Every read
that resolves takes \texttt{as\_of} and \texttt{as\_known}, each defaulting to
the server clock.}
\label{tab:surface}
\end{table}

The extraction pair reads only relations a document states in the form
\texttt{relation:: value}, with \texttt{[[key]]} marking a value that names an
entity, and files them through the same admission as \op{assert}; it infers
nothing from prose. Full-text search is an SQLite FTS5 index with external
content over the same rows, so it holds no second copy of the text.
